\documentclass[11pt]{article}
\usepackage[margin=1in]{geometry}
\usepackage{amsmath,amssymb,amsthm}
\usepackage{hyperref}

\newtheorem{theorem}{Theorem}
\newtheorem{proposition}{Proposition}
\newtheorem{lemma}{Lemma}
\newtheorem{conjecture}{Conjecture}
\newtheorem{remark}{Remark}
\newcommand{\mc}{\operatorname{MaxCut}}
\newcommand{\dd}{d}
\newcommand{\E}{\mathbb{E}}

\title{Exchange Identities and Computational Diagnostics for Triangle-Free Graph Bipartization}
\author{Jinji Li}
\date{Version 2.0 draft -- September 2026}

\begin{document}
\maketitle

\begin{abstract}
For a graph $G$, let
\[
\dd(G)=|E(G)|-\mc(G)
\]
be the minimum number of edges whose deletion makes $G$ bipartite.
A classical conjecture of Erd\H{o}s asserts that every triangle-free graph
on $N$ vertices satisfies $\dd(G)\le N^2/25$; for $N=5k$ this is
$\dd(G)\le k^2$.

We study a five-vertex induction strategy whose target is
\[
\exists X\subseteq V(G),\ |X|=5:\qquad
\dd(G)-\dd(G-X)\le 2k-1.
\]
The paper develops an exact coordinated-flip identity, a corresponding
selection inequality, elementary edge/star/biclique interaction
corollaries, and reproducible finite certificates for reduced extension
models.  We also record exact computational diagnostics showing that
stronger fixed-five-set statements can fail even when the existential
induction target is easily repaired by changing the five-set.

A correction to an earlier draft is important.  If
\[
\ell(X)=2\dd(G[X])+e(X,G-X),
\]
then the empty flip already proves the induction target whenever
$\ell(X)\le 4k-1$.  Thus at $k=3$ the automatic range extends through
$L(G)=11$, where $L(G)=\min_{|X|=5}\ell(X)$.  The previously reported
$L=11$ exchange audit therefore remains a valid structural computation,
but is not needed to prove the induction inequality in that regime.

We place the project in the established language of frustration-critical
signed graphs and weakly bipartite signed graphs.  The full Erd\H{o}s
conjecture remains open; this manuscript is a structural and computational
investigation of one proof architecture, not a proof of the conjecture.
\end{abstract}

\section{Introduction}

For a finite simple graph $G$, define
\[
\dd(G)=\min\{|F|:F\subseteq E(G),\ G-F\text{ is bipartite}\}
      =|E(G)|-\mc(G).
\]
Erd\H{o}s conjectured that every triangle-free graph on $N$ vertices
satisfies
\[
\dd(G)\le \frac{N^2}{25}.
\]
The balanced blow-up of $C_5$ shows that the constant $1/25$ would be
sharp. Balogh, Clemen and Lidick\'y proved the global bound $N^2/23.5$
and proved the sharp $N^2/25$ bound in two density ranges
\cite{BCL}.  A recent computer-assisted preprint of Ferudun proves the
sharp value for every multiple of five up to $N=200$ \cite{Ferudun}.
Consequently, the small finite cases appearing below are used as
validation and structure-finding devices, not as new extremal evaluations.

The induction target studied here is the following.

\begin{conjecture}[Five-set induction target]
For every triangle-free graph $G$ on $5k$ vertices, there exists a
five-set $X\subseteq V(G)$ such that
\[
\dd(G)-\dd(G-X)\le 2k-1.
\]
\end{conjecture}

If this statement holds for all $k$, induction from $k=1$ gives
\[
\dd(G)\le 1+3+\cdots+(2k-1)=k^2.
\]
Thus the problem becomes one of selecting a good five-set and extending an
optimal coloring of the remaining graph.

This paper has two goals.  First, we develop exact identities that separate
boundary imbalance from internal recoloring gain.  Second, we use finite
certificates to diagnose which strengthenings of the induction are
plausible and which are false.

\section{Literature context and scope}

The quantity $\dd(G)$ is the edge odd-cycle-transversal number, equivalently
the MaxCut defect.  In the language of signed graphs it is a special case
of the frustration index.  Cappello and Steffen study $k$-critical signed
graphs, where deleting any edge lowers the frustration index
\cite{CappelloSteffen}.  We therefore do not claim the general notion of
edge-criticality as new.

Likewise, the odd-circuit-cover polyhedron and its integrality theory are
classical.  Guenin's theorem characterizes weakly bipartite signed graphs
by exclusion of an odd-$K_5$ minor; Schrijver gave a short proof, and
Geelen--Guenin developed related packing consequences
\cite{Schrijver,Guenin}.  These results are used here only as background
for possible obstruction viewpoints.

The contribution claimed in this manuscript is narrower: the specific
five-set induction framework, the coordinated-flip selection identity and
its use as a diagnostic tool, exact reduced-state certificates, and
reproducible computations that distinguish fixed-set failure from
existential five-set success.

\section{Exact finite ground truth}

All exact computations use
\[
\dd(G)=|E(G)|-\mc(G).
\]
As a regression computation, all $12{,}172$ triangle-free isomorphism
classes on ten vertices were checked exactly, with
\[
\max \dd(G)=4,
\]
attained by the balanced $C_5$ blow-up $B_2$.  In view of
\cite{Ferudun}, this computation is not a novelty claim.

\section{Balanced and mixed extension certificates}

The first reduction studies finite extension models organized around the
five-part structure of balanced $C_5$ blow-ups.

\subsection{Balanced-extension certificate}

For $t=1,\ldots,25$, the exact C++ checker
\texttt{scripts/check\_balanced\_extension.cpp} enumerates the reduced
balanced-extension states.  In every batch the reported maximum gap is
zero.

\begin{proposition}[Balanced-extension certificate]
For every balanced-extension state represented by the reduction at
$t=1,\ldots,25$, the target inequality checked by the exact program holds.
No violating state occurs in the enumerated reduced state space.
\end{proposition}

\subsection{Exact mixed $t=2$ closure}

The mixed $t=2$ model contains $4{,}013{,}521$ maximal-pair states, of
which $2{,}768{,}154$ are genuinely mixed.  Exact MaxCut evaluation gives
maximum deletion distance $9$ over the full state space and maximum $8$
among the genuinely mixed states.  A maximal-completion monotonicity
argument reduces arbitrary represented legal states to maximal
triangle-free completions.

\begin{proposition}[Mixed $t=2$ certificate]
Within the exact mixed $t=2$ extension model, all represented legal states
are closed by maximal triangle-free completion.  Among maximal states the
largest deletion distance is $9$, and among genuinely mixed states it is
$8$.
\end{proposition}

These certificates motivated allowing the coloring of the core graph to
change in a coordinated way.

\section{Coordinated flips}
\label{sec:flips}

Fix a five-set $X\subseteq V(G)$ and put $H=G-X$.  Let $c$ be an optimal
2-coloring of $H$.  For an assignment $a$ on $X$, let $E_c(a)$ denote the
number of monochromatic edges contributed by $G[X]$ and $E(X,H)$ before
modifying the core coloring.

For $v\in V(H)$ define
\[
w_a(v)=
\#\{\text{monochromatic edges from $v$ to $X$}\}
-
\#\{\text{bichromatic edges from $v$ to $X$}\},
\]
and
\[
\sigma_c(v)=
\#\{\text{bichromatic core edges incident with $v$}\}
-
\#\{\text{monochromatic core edges incident with $v$}\}.
\]
For $S\subseteq V(H)$, let $m_c(S)$ be the number of monochromatic edges
of $H[S]$ under $c$.

\begin{lemma}[Coordinated-flip identity]
If all colors on $S$ are flipped, then
\[
g_c(a,S)=
\sum_{v\in S}\bigl(w_a(v)-\sigma_c(v)\bigr)
+2e(H[S])-4m_c(S)
\]
satisfies
\[
 b_G(c^S\cup a)-\dd(H)=E_c(a)-g_c(a,S),
\]
where $b_G(\cdot)$ is the number of monochromatic edges under the displayed
coloring.
\end{lemma}

\begin{proof}
Each boundary edge incident with a flipped vertex changes according to the
signed contribution $w_a(v)$.  Each core edge with exactly one endpoint
flipped changes according to $-\sigma_c(v)$ after summing over vertices.
Internal edges of $H[S]$ are counted twice by the vertex sum and require
the correction $2e(H[S])-4m_c(S)$.  Collecting these contributions gives
the formula.
\end{proof}

The implemented admissible flip family consists of independent sets and
sets inducing complete bipartite subgraphs, together with their
complements.

Let $\mu_X$ be the uniform distribution on optimal 2-colorings of $G[X]$,
counting color reversals separately.  For fixed optimal $c$,
\[
\E_{a\sim\mu_X}E_c(a)
=\dd(G[X])+\frac{e(X,H)}2.
\]
Define
\[
\Phi_X=
\max_{c\text{ optimal on }H}
\E_{a\sim\mu_X}\max_{S\text{ admissible}}g_c(a,S).
\]

\begin{proposition}[Coordinated-flip selection bound]
For every five-set $X$,
\[
\dd(G)-\dd(G-X)
\le
\left\lfloor
\dd(G[X])+\frac{e(X,G-X)}2-\Phi_X
\right\rfloor.
\]
\end{proposition}

\begin{proof}
For each assignment $a$, choose an admissible flip maximizing the gain.
The resulting coloring gives an upper bound on $\dd(G)$.  Average over
$a\sim\mu_X$ and then maximize over the optimal core coloring $c$.
Integrality gives the floor.
\end{proof}

\subsection{Bichromatic biclique corollary}

Write
\[
x_a(v)=w_a(v)-\sigma_c(v).
\]
If $S=A\cup B$ induces a complete bipartite subgraph of $H$ and every
internal edge is bichromatic under $c$, then $m_c(S)=0$ and
$e(H[S])=|A||B|$.  Hence the flip identity immediately gives
\[
\boxed{
 g_c(a,S)=\sum_{v\in S}x_a(v)+2|A||B|.
}
\]
The pair and star formulas are the cases $K_{1,1}$ and $K_{1,t}$.
These are exact bookkeeping corollaries rather than separate extremal
results, but they are useful for identifying interaction gain that is
invisible in singleton flips.

\section{The exact empty-flip threshold}

Define
\[
\ell(X)=2\dd(G[X])+e(X,G-X),
\qquad
L(G)=\min_{|X|=5}\ell(X).
\]
Setting $S=\varnothing$ gives $\Phi_X\ge0$ and therefore
\[
\dd(G)-\dd(G-X)
\le
\left\lfloor\frac{\ell(X)}2\right\rfloor.
\]

\begin{proposition}[Automatic five-set threshold]
For a triangle-free graph on $5k$ vertices, if a five-set $X$ satisfies
\[
\ell(X)\le 4k-1,
\]
then
\[
\dd(G)-\dd(G-X)\le2k-1.
\]
In particular, Candidate A is automatic whenever $L(G)\le4k-1$.
\end{proposition}

\begin{remark}[Correction to Version 1]
At $k=3$, the automatic threshold is $L(G)\le11$, not $L(G)\le10$.
Thus global $L(G)=11$ is not a genuinely hard Candidate-A regime.  The
$L=11$ computations below remain valid as experiments on exchange
structure, but they are not required to prove the five-set inequality.
\end{remark}

\section{Fixed-set failure and exchange diagnostics}

Early experiments suggested stronger fixed-five-set gain statements.  A
direct actual-core analysis produced realizable residual configurations
for which these stronger gain thresholds fail.  Importantly, failure of a
strong fixed-$X$ statement does not imply failure of the existential
five-set induction target.

For $22$ residual witnesses, an exact audit over all
$\binom{15}{5}=3003$ five-sets gives
\[
\max q_{\min}=4,
\qquad
q(Y)=\dd(G)-\dd(G-Y),
\]
and restricting to 3-for-3 exchanges relative to the original set gives
restricted minima distributed as
\[
\begin{array}{c|rrrr}
q_{\min} & 1 & 2 & 3 & 4\\ \hline
\# & 1 & 10 & 8 & 3.
\end{array}
\]
Thus every one of these explicit fixed-set obstructions is repaired by
changing the five-set.

\section{The saved global-$L=11$ 3-for-3 audit}

A 3-for-3 exchange of a recorded five-set $X$ is a five-set $Y$ satisfying
$|X\cap Y|=2$.  On fifteen vertices there are
\[
\binom52\binom{10}{3}=1200
\]
such exchanges.

The saved corpus contains $3{,}436$ canonical records with global
$L(G)=11$.  Relative to the distinguished recorded five-set, every one of
these records admits a 3-for-3 exchange $Y$ with
\[
\dd(G)-\dd(G-Y)\le5.
\]
Among them, $762$ have $\dd(G)>5$ and require a nontrivial search; all
$762$ pass.

\begin{remark}[Interpretation after the threshold correction]
Because $L(G)=11$ already lies in the automatic range at $k=3$, this audit
is not a proof certificate needed for Candidate A.  Its value is
structural: it demonstrates, on a precisely scoped saved corpus, that
obstructions to stronger fixed-set statements can be repaired by a small
exchange of the deletion set.
\end{remark}

The scope audit independently matches the $3{,}436$ global-$L=11$ IDs and
canonical graph6 strings against the exchange file.  The statement is
exhaustive only for this saved global-$L=11$ corpus; no claim is made here
that the parent saved corpus is the full universe of triangle-free graphs
on fifteen vertices.

\section{Criticality and odd-cycle-cover language}

The edge-bipartization number is the minimum size of an edge set meeting
every odd cycle.  The following two elementary facts are useful when
considering a hypothetical edge-minimal counterexample.

\begin{lemma}[Critical edge and minimum transversal]
For an edge $e$,
\[
\dd(G-e)=\dd(G)-1
\]
if and only if $e$ belongs to some minimum odd-cycle edge transversal of
$G$.
\end{lemma}

\begin{proof}
If a minimum transversal $T$ contains $e$, then $T\setminus\{e\}$
bipartizes $G-e$, so $\dd(G-e)\le\dd(G)-1$; deleting one edge can reduce
$\dd$ by at most one, hence equality holds.  Conversely, if
$\dd(G-e)=\dd(G)-1$, a minimum transversal of $G-e$, together with $e$,
is a minimum transversal of $G$ containing $e$.
\end{proof}

\begin{lemma}[Private odd-cycle witness]
Let $T$ be an inclusion-minimal odd-cycle edge transversal.  For every
$e\in T$ there exists an odd cycle $C_e$ with
\[
C_e\cap T=\{e\}.
\]
If $G$ is triangle-free, then $|C_e|\ge5$.
\end{lemma}

\begin{proof}
If no such cycle existed, then every odd cycle meeting $e$ would also meet
$T\setminus\{e\}$, so $T\setminus\{e\}$ would still meet every odd cycle,
contradicting inclusion-minimality.
\end{proof}

These observations belong naturally to the existing theory of
frustration-critical signed graphs \cite{CappelloSteffen}.  They are
included to connect the five-set induction project with that language, not
as novelty claims.

Weakly bipartite signed graphs provide another relevant viewpoint.  The
odd-circuit-cover polyhedron is integral exactly in the odd-$K_5$-minor-free
class by Guenin's theorem \cite{Schrijver,Guenin}.  Any future LP-based
attack on the Erd\H{o}s conjecture must therefore distinguish genuinely new
triangle-free consequences from this established signed-graph theory.

\section{Reproducibility}

The principal release artifacts are:
\begin{itemize}
\item exact balanced- and mixed-extension checkers;
\item \texttt{results/hard\_regime\_L\_n15/L\_records.tsv};
\item \texttt{results/hard\_regime\_L\_n15/L11\_exchange.jsonl};
\item \texttt{scripts/build\_L11\_3for3\_input.py};
\item \texttt{scripts/check\_L11\_3for3\_all.cpp};
\item deterministic audit outputs and SHA-256 checksums.
\end{itemize}
The computational claims are intended to be reproducible from explicit
finite inputs.  We distinguish throughout between symbolic identities,
finite exact certificates, and conjectural structural statements.

\section{Discussion}

The main lesson of the computation is negative as well as positive.
Strong fixed-five-set statements can fail for realizable configurations,
while the existential induction target remains true after changing the
five-set.  This suggests that a successful proof, if it follows this
architecture, should prioritize selection over increasingly elaborate
certification of a distinguished set.

The corrected floor threshold further simplifies the picture: at $k=3$,
all global-$L=11$ records are already automatic for Candidate A.  The
exchange audit should therefore be viewed as a diagnostic experiment on
selection geometry rather than as evidence needed to close that regime.

The critical signed-graph and weakly-bipartite viewpoints provide a useful
external check on novelty.  Edge-criticality, private odd-cycle witnesses,
and odd-$K_5$ integrality obstructions are part of established theory.  A
future theorem paper should therefore contribute a genuinely new
triangle-free structural statement connecting these tools to the sharp
$N^2/25$ bound.

\section*{Status and non-claims}

The Erd\H{o}s triangle-free bipartization conjecture remains open in
general.  This manuscript does not claim to solve it.  The finite
computations at $n=10$ and $n=15$ are not claimed as new extremal
evaluations.  The global-$L=11$ 3-for-3 audit is a structural computation,
not a necessary proof of Candidate A after the corrected threshold.
The manuscript's current contribution is a reproducible five-set induction
framework, exact flip identities, reduced-state certificates, and explicit
diagnostics of false strengthenings.

\begin{thebibliography}{99}
\bibitem{BCL}
J. Balogh, F. C. Clemen, and B. Lidick\'y,
\emph{Max Cuts in Triangle-Free Graphs},
Extended Abstracts EuroComb 2021, Trends in Mathematics, Springer, 2021;
arXiv:2103.14179.

\bibitem{Ferudun}
A. Ferudun,
\emph{The Erd\H{o}s $n^2/25$ max-cut conjecture for small multiples of five,
via a per-root-MaxCut envelope and blow-up integrality},
arXiv:2606.28041, 2026.

\bibitem{CappelloSteffen}
C. Cappello and E. Steffen,
\emph{Frustration-critical signed graphs},
Discrete Applied Mathematics 322 (2022), 183--193;
arXiv:2112.02664.

\bibitem{Schrijver}
A. Schrijver,
\emph{A short proof of Guenin's characterization of weakly bipartite graphs},
Journal of Combinatorial Theory, Series B 85 (2002), 79--83.

\bibitem{Guenin}
J. F. Geelen and B. Guenin,
\emph{Packing Odd Circuits in Eulerian Graphs},
Journal of Combinatorial Theory, Series B 86 (2002), 280--295.
\end{thebibliography}

\end{document}
